% ────────────────────────────────────────────────────────────────────────
%  LaTeX Source for Methodology and Results & Discussion Sections
%  Format: IEEE Conference Template (IEEEtran)
%  Includes CPU (SW/HW 1,2,4 cores) and GPU (CUDA v1 & v2 / 4,16,64,256 KB)
% ────────────────────────────────────────────────────────────────────────

\section{Experimental Methodology}
\label{sec:methodology}

To evaluate the performance and energy efficiency of dedicated hardware cryptographic extensions against traditional software implementations, we deployed a heterogeneous benchmark environment consisting of two distinct compute platforms:
\begin{itemize}
    \item \textbf{CPU Platform:} An Intel Core i7-1255U processor (12th Gen Alder Lake) featuring 2 Performance cores (P-cores) and 8 Efficient cores (E-cores). All CPU workloads were pinned exclusively to the E-cores (fixed base frequency of 3.5~GHz, Hyper-Threading disabled) using the Linux \texttt{taskset} utility (CPUs 4--7), eliminating interference from heterogeneous core asymmetry and dynamic frequency scaling.
    \item \textbf{GPU Platform:} An NVIDIA GeForce RTX 2050 mobile GPU (GA107M, Ampere architecture, \texttt{sm\_86}), featuring 16 Streaming Multiprocessors (SMs), 1,536 threads/SM, 65,536 registers/SM, 4~GB of GDDR6 VRAM at 94.4~GB/s peak bandwidth.
\end{itemize}

We evaluated three distinct SHA-256 implementation paradigms:
\begin{enumerate}
    \item \textbf{CPU Software (SW):} A standard scalar C implementation compiled with GCC using \texttt{-mno-sha} and \texttt{-march=x86-64} flags to explicitly disable hardware-accelerated SHA instructions, confining computation to the general-purpose ALU pipeline.
    \item \textbf{CPU Hardware-Accelerated (HW / SHA-NI):} An OpenSSL 3.5.5 EVP-based implementation that performs runtime CPUID dispatch to dedicated Intel SHA-NI instructions (\texttt{SHA256RNDS2}, \texttt{SHA256MSG1}, \texttt{SHA256MSG2}).
    \item \textbf{GPU Hardware-Accelerated (CUDA):} A parallel CUDA C/C++ implementation partitioning the workload into independent fixed-size chunks, with one CUDA thread computing the SHA-256 digest of each chunk in parallel. Two kernel versions were evaluated:
    \begin{itemize}
        \item \textbf{v1 (baseline):} Standard 64-register message schedule (\texttt{w[64]}), global memory for the SHA-256 K constants.
        \item \textbf{v2 (high-utilization):} Eight targeted optimizations to maximize SM occupancy and instruction throughput (Section~\ref{sec:gpu_opt}).
    \end{itemize}
\end{enumerate}

The workload consisted of hashing a 1.0~GB binary payload pre-warmed into the OS page cache to eliminate disk I/O latency. CPU energy was measured via Intel RAPL (\texttt{/sys/class/powercap/intel-rapl:0/energy\_uj}). GPU power draw was sampled at 100~ms intervals using \texttt{nvidia-smi} polling during each run to compute average watt draw and total Joules consumed. Five iterations were executed per configuration to obtain statistically stable means.

\subsection{GPU Kernel Optimizations}
\label{sec:gpu_opt}

The baseline GPU kernel (v1) suffered from low SM occupancy (~33\%) due to high register pressure: the standard SHA-256 message schedule expansion requires a full \texttt{w[64]} array (256 bytes = 64 registers per thread), combined with 8 state registers and 10 working variables, totalling approximately 100 live registers per thread. At 256 threads/block, this allows only $\lfloor 65536 / (100 \times 256) \rfloor = 2$ concurrent blocks per SM, or 512 of the 1,536 possible concurrent threads.

Version v2 applies the following optimizations to reduce register pressure and maximize hardware utilization:
\begin{enumerate}
    \item \textbf{Sliding-window message schedule} (\texttt{w[16]}): The standard 64-word expansion is replaced by a 16-word circular buffer. The recurrence $w_i = \sigma_1(w_{i-2}) + w_{i-7} + \sigma_0(w_{i-15}) + w_{i-16}$ is computed in-place by overwriting index $i \bmod 16$. This reduces message schedule registers from 64 to 16 (saving 192 bytes per thread).
    \item \textbf{\texttt{\_\_constant\_\_} memory for K}: The 256-byte constant table is placed in CUDA constant memory, enabling a single hardware broadcast transaction per warp instead of 32 independent L2 cache reads.
    \item \textbf{\texttt{\#pragma unroll 64}} on compression rounds: The compiler emits straight-line PTX for all 64 rounds, maximizing instruction-level parallelism and eliminating branch misprediction overhead.
    \item \textbf{\texttt{\_\_byte\_perm()} big-endian word load}: A single PTX \texttt{sel.b32} instruction performs the byte-swap required for SHA-256 big-endian word loading, replacing four shift-and-OR operations.
    \item \textbf{Inline padding block}: The SHA-256 length-padding block is materialized directly as 16 known-value registers rather than a 64-byte local memory buffer, eliminating stack spill.
    \item \textbf{Grid-stride loop}: Instead of launching one thread per chunk, a fixed thread pool of $16~\text{SMs} \times 5~\text{blocks/SM} \times 256~\text{threads/block} = 20{,}480$ threads covers all chunks via strided iteration, guaranteeing 100\% SM saturation from first to last chunk.
    \item \textbf{Pinned host memory} (\texttt{cudaMallocHost}): Page-locked allocation enables asynchronous DMA transfers, increasing effective PCIe bandwidth by up to 2$\times$.
    \item \textbf{Triple CUDA streams}: The file is divided into three equal segments processed concurrently on independent streams, overlapping host-to-device transfer, kernel execution, and device-to-host transfer simultaneously using separate DMA and compute engines.
\end{enumerate}

The combined effect raises SM occupancy from 33\% to approximately 83\% (5 resident blocks/SM at $\leq$48 live registers/thread, enforced via \texttt{--maxrregcount=48}).

\section{Results and Discussion}
\label{sec:results}

Tables~\ref{tab:cpu_results} and~\ref{tab:gpu_results} report the complete benchmark results across all evaluated configurations.

% ── CPU Results Table ─────────────────────────────────────────────────────────
\begin{table}[htbp]
\centering
\caption{SHA-256 CPU Benchmark (Intel Core i7-1255U E-cores @ 3.5~GHz)}
\label{tab:cpu_results}
\begin{tabular}{lcccc}
\hline\hline
\textbf{Impl.} & \textbf{Threads} & \textbf{Throughput (MB/s)} & \textbf{Speedup} & \textbf{Energy (J/GB)} \\
\hline
\textbf{SW (C)} & 1 & $189.73 \pm 1.20$ & $1.00\times$ & 69.76 \\
                & 2 & $372.18 \pm 0.90$ & $1.96\times$ & 44.23 \\
                & 4 & $735.65 \pm 1.80$ & $3.88\times$ & 32.89 \\
\hline
\textbf{HW (SHA-NI)} & 1 & $1,763.18 \pm 4.10$ & $9.29\times$ & 7.00 \\
                     & 2 & $3,315.00 \pm 22.30$ & $17.47\times$ & 5.37 \\
                     & 4 & $5,126.07 \pm 320.10$ & $27.02\times$ & 5.17 \\
\hline\hline
\end{tabular}
\end{table}

% ── GPU Results Table ─────────────────────────────────────────────────────────
\begin{table}[htbp]
\centering
\caption{SHA-256 GPU Benchmark (NVIDIA RTX 2050, sm\_86, 1.0~GB payload)}
\label{tab:gpu_results}
\begin{tabular}{llcccc}
\hline\hline
\textbf{Version} & \textbf{Chunk} & \textbf{Kernel (GB/s)} & \textbf{E2E (GB/s)} & \textbf{Speedup\textsuperscript{†}} & \textbf{Energy (J/GB)} \\
\hline
\textbf{v1} & 4~KB   & $3.57 \pm 0.02$ & $1.92 \pm 0.01$ & $18.8\times$ & 4.74 \\
(Baseline)  & 16~KB  & $3.24 \pm 0.04$ & $1.77 \pm 0.06$ & $17.1\times$ & 5.73 \\
            & 64~KB  & $2.81 \pm 0.01$ & $1.67 \pm 0.04$ & $16.2\times$ & 6.30 \\
            & 256~KB & $7.41 \pm 0.05$ & $2.60 \pm 0.02$ & $39.1\times$ & 3.77 \\
\hline
\textbf{v2} & 4~KB   & $11.31 \pm 0.38$ & $4.42 \pm 0.05$ & $59.7\times$ & \textbf{2.04} \\
(High-Util) & 16~KB  & $10.75 \pm 0.12$ & $4.51 \pm 0.13$ & $57.0\times$ & 2.25 \\
            & \textbf{64~KB}  & $\mathbf{16.81 \pm 0.48}$ & $4.61 \pm 0.18$ & $\mathbf{88.6\times}$ & 2.16 \\
            & 256~KB & $11.07 \pm 0.10$ & $3.71 \pm 0.01$ & $58.4\times$ & 2.77 \\
\hline\hline
\multicolumn{6}{l}{\textsuperscript{†}Kernel speedup vs. CPU SW single-thread baseline (189.73~MB/s).} \\
\end{tabular}
\end{table}

\subsection{CPU Throughput and SHA-NI Hardware Acceleration}

In single-threaded execution, Intel SHA-NI (\texttt{HW 1T}) achieves 1,763.18~MB/s, a \textbf{9.29$\times$} speedup over the software baseline. When scaled to 4 E-cores, SHA-NI reaches 5,126.07~MB/s but exhibits degraded scaling efficiency (72.7\%), while the pure software implementation scales near-linearly at 96.9\% efficiency. This divergence is explained by the memory bandwidth limit: at 5.1~GB/s of aggregate cryptographic throughput, four SHA-NI pipelines saturate the dual-channel LPDDR4x memory subsystem, shifting the bottleneck from compute to memory bandwidth.

\subsection{GPU v1 Baseline: Compute vs. PCIe Bottleneck}

The GPU v1 baseline demonstrates strong single-SM computation (7.41~GB/s kernel at 256~KB chunks), a \textbf{39.1$\times$} speedup over the CPU software baseline. However, the end-to-end throughput collapses to 2.60~GB/s due to PCIe host-to-device and device-to-host data transfers. This confirms that for GPU cryptographic workloads, the PCIe interconnect is the primary system bottleneck, limiting real-world throughput to approximately 35\% of on-device compute capacity.

\subsection{GPU v2: Impact of High-Utilization Optimizations}

Version v2 demonstrates dramatic improvements across all chunk sizes. The peak kernel throughput of \textbf{16.81~GB/s} is achieved at 64~KB chunks, representing a \textbf{+499\%} improvement over the v1 baseline at the same configuration (2.81~GB/s), and a \textbf{+127\%} improvement over v1's best configuration (7.41~GB/s at 256~KB).

The dominant contributor to this improvement is the reduction of register pressure through the sliding-window message schedule. By reducing live registers from approximately 100 to 48 per thread, the number of resident blocks per SM increases from 2 to 5, raising SM occupancy from 33\% to 83\%. Combined with the unrolled compression rounds and constant-memory K table broadcast, this substantially increases the instruction throughput and latency-hiding capability of each SM.

The optimal chunk size shifts from 256~KB in v1 to \textbf{64~KB} in v2. In v1, large chunks compensate for the low parallelism caused by restricted occupancy. In v2, the grid-stride kernel ensures full SM saturation at all chunk sizes; the 64~KB configuration provides sufficient thread-level parallelism (16,384 concurrent chunks) while maintaining cache-friendly sequential memory access patterns within each thread.

\subsection{Energy Efficiency Across Paradigms}

Table~\ref{tab:energy_comparison} presents a cross-platform energy efficiency comparison.

\begin{table}[htbp]
\centering
\caption{Energy Efficiency Comparison (J/GB, 1.0~GB payload)}
\label{tab:energy_comparison}
\begin{tabular}{lcc}
\hline\hline
\textbf{Platform / Config} & \textbf{Throughput} & \textbf{Energy (J/GB)} \\
\hline
CPU SW, 1T              & 189.7~MB/s  & 69.76 \\
CPU SW, 4T              & 735.7~MB/s  & 32.89 \\
CPU SHA-NI, 1T          & 1,763~MB/s  & 7.00  \\
CPU SHA-NI, 4T          & 5,126~MB/s  & 5.17  \\
GPU v1, 256~KB (E2E)    & 2,662~MB/s  & 3.77  \\
GPU v2, 64~KB (E2E)     & 4,717~MB/s  & 2.16  \\
\textbf{GPU v2, 4~KB (kernel)} & \textbf{11,571~MB/s} & \textbf{1.90}  \\
\hline\hline
\end{tabular}
\end{table}

The GPU v2 kernel-only metric (1.90~J/GB) represents a \textbf{36.7$\times$} improvement in energy efficiency over the CPU software baseline. Even considering the full end-to-end pipeline including PCIe transfers (2.16~J/GB at 64~KB), the GPU remains \textbf{2.4$\times$} more energy efficient than the 4-core CPU SHA-NI configuration (5.17~J/GB) while delivering higher total throughput. This validates the GPU as the optimal platform for high-volume, batch-mode SHA-256 workloads from both performance and energy perspectives.
